\section{Symbols, Abbreviations and Acronyms}

\renewcommand{\arraystretch}{1.2}
\begin{tabular}{l l} 
  \toprule		
  \textbf{symbol} & \textbf{description}\\
  \midrule 
  T & Test\\
  FR & Functional Requirement (from SRS)\\
  NFR & Nonfunctional Requirement (from SRS)\\
  VnV & Verification and Validation\\
  API & Application Programming Interface\\
  DB & Database\\
  M1 & Authentication \& Verification Module\\
  M2 & User Profile Module\\
  M3 & Route \& Trip Module\\
  M4 & Matching Module\\
  M5 & Scheduling Module\\
  M6 & Notification Module\\
  M7 & Rating \& Feedback Module\\
  M8 & Safety \& Reporting Module\\
  M9 & Payment \& Cost Estimation Module\\
  M10 & Pricing Module\\
  M11 & Database Module\\
  \bottomrule
\end{tabular}\\

\newpage

\tableofcontents

\listoftables %if appropriate

\listoffigures %if appropriate

\newpage

\pagenumbering{arabic}

This document presents the Verification and Validation (VnV) activities and results for Hitchly. The symbols and abbreviations listed above are used consistently throughout the report. In particular, IDs of the form \texttt{test-FR*}, \texttt{test-R*}, and \texttt{test-U*} denote system-level tests traced to functional and nonfunctional requirements (Sections~3, 4, and 9), whereas IDs of the form \texttt{test-ut-*} denote code-level unit tests traced to modules (Sections~6 and 10). Module identifiers M1--M11 are used to refer to the design modules defined in the MG/MIS and summarized in this section.
